\documentclass[12pt]{article}
\usepackage[spanish]{babel}
\usepackage{natbib}
\usepackage{url}
\usepackage[utf8x]{inputenc}
\usepackage[mediumspace,mediumqspace,squaren]{SIunits}
\usepackage{amsmath}
\usepackage{graphicx}
\graphicspath{{images/}}
\usepackage{parskip}
\usepackage{fancyhdr}
\usepackage{vmargin}
\usepackage{enumerate} % enumerados
\usepackage{diagbox}
\usepackage{booktabs}
\usepackage{float} %para usar [H]


\setmarginsrb{3 cm}{2.5 cm}{3 cm}{2.5 cm}{1 cm}{1.5 cm}{1 cm}{1.5 cm}

\title{Practico de Laboratorio N°1: Adaptación en lineas de transmisión }								% Title
\author{}								% Author
\date{2015}											% Date

\makeatletter
\let\thetitle\@title
\let\theauthor\@author
\let\thedate\@date
\makeatother

\pagestyle{fancy}
\fancyhf{}
\rhead{\theauthor}
\lhead{\thetitle}
\cfoot{\thepage}

\begin{document}

%%%%%%%%%%%%%%%%%%%%%%%%%%%%%%%%%%%%%%%%%%%%%%%%%%%%%%%%%%%%%%%%%%%%%%%%%%%%%%%%%%%%%%%%%

\begin{titlepage}
	\centering
    \vspace*{0.5 cm}
     \includegraphics[scale = 0.8]{Logo_UTN.png}\\[1.0 cm]	% University Logo
    \textsc{\LARGE Universidad Tecnológica Nacional}\\[2.0 cm]	% University Name
	\textsc{\Large Curso 3R2}\\[0.5 cm]				% Course Code
	\textsc{\large Medios de Enlace}\\[0 cm]				% Course Name
	\textrm{\large (Ing. Contreras, Luis Candelario)}\\[0.5 cm]
    \rule{\linewidth}{0.2 mm} \\[0.4 cm]
	{ \huge \bfseries \thetitle}\\
	\rule{\linewidth}{0.2 mm} \\[1 cm]
	
	\begin{minipage}{0.4\textwidth}
		\begin{flushleft} \large
			\emph{Autores:}\\
			Bianchini, Bruno\\Fichetti P. Tomás\\Guizzo, José\\Bosse, Esteban A.
			\end{flushleft}
			\end{minipage}~
			\begin{minipage}{0.4\textwidth}
			\begin{flushright} \large
			\emph{Legajo:} \\
			 63070\\62281\\62165\\62930									% Your Student Number
		\end{flushright}
	\end{minipage}\\[2 cm]
	
	{\large \thedate}\\[2 cm]
 
	\vfill
	
\end{titlepage}

%%%%%%%%%%%%%%%%%%%%%%%%%%%%%%%%%%%%%%%%%%%%%%%%%%%%%%%%%%%%%%%%%%%%%%%%%%%%%%%%%%%%%%%%%

\tableofcontents
\pagebreak

%%%%%%%%%%%%%%%%%%%%%%%%%%%%%%%%%%%%%%%%%%%%%%%%%%%%%%%%%%%%%%%%%%%%%%%%%%%%%%%%%%%%%%%%%
\section{Objetivo General}

\begin{itemize}
 \item Analizar la respuesta a un escalón de tensión de una línea de trasmisión coaxial de impedancia característica \textbf{Zo} terminada en distintos valores de impedancia de carga ZL.
 \item Analizar cualitativamente la manera de determinar la impedancia característica (\textbf{Zo}) de una línea de transmisión por observación de la disminución y/o anulación de la distorsión de la señal en la carga.
 \item Medir el tiempo de retardo (\textbf{tr}) en una línea de trasmisión.
 \item Medir la atenuación (\textbf{At}) en decibeles (\textbf{db}) de una línea de trasmisión.
 \item Medir la constante de propagación (\textbf{K = Vp / C}) en una línea de trasmisión.
 \end{itemize} 

\section{Conclusiones}
Cada grupo de trabajo (no más de 3 alumnos por grupo) deberá realizar un informe, el que deberá ser presentado a su respectivo Profesor de Cátedra antes de la finalización del año lectivo 2015, como condición adicional a las existentes para lograr la regularización de la materia.
El informe deberá ser presentado en hojas formato A4 a faz simple y deberá tener:

\begin{itemize}
 \item La presente guía de TP llenada con los valores requeridos.
 \item Mediciones y cálculos realizados y gráficos observados.
 \item Expresar brevemente el marco teórico que se relaciona con este trabajo práctico.
 \item Indicar la experiencia adquirida en la realización del trabajo práctico en el Laboratorio de Comunicaciones.
 \item Realizar sugerencias (si las hubiere) para futuros trabajos prácticos.
 \end{itemize} 



\section{Materiales Requeridos}
Los elementos utilizados en siguiente práctico son:
\begin{itemize}


 \item Generador de onda cuadrada (f = 100 KHz; V= 3 Vpp)
 \item Osciloscopio de doble trazo.
 \item Multímetro.
 \item Cable coaxial de 75 $\Omega$  (305 metros)
 \item Kit de adaptación en líneas de transmisión del LdC.
\end{itemize}


\section{Actividades}

\subsection{Actividad N°1}
 En base al circuito propuesto, conectar una resistencia de carga ZL= 75 $\Omega$ en el punto \textbf{B}; observar las formas de onda en los puntos \textbf{A} y \textbf{B} (la línea de transmisión está adaptada: Zo = ZL = 75 $\Omega$). 

 \begin{center}

  TIME/DIV: 2 $\micro$S

  Ch1:2 V/div \hspace{0.5cm} $V_{pp}$= 2 x 2  \hspace{0.5cm}   $V_{pp}$=4 V
  
  Ch2: 2 V/div \hspace{0.5cm} $V_{pp}$= 1.6 x 2 \hspace{0.5cm}   $V_{pp}$=3.2 V
\end{center}
                                                  
 
\subsection{Actividad N°2}
Graficar la forma de onda de los canales Ch1 y Ch2.

\begin{figure}[hbt]
\centering
\includegraphics[scale = 0.3]{imagen7.jpg}\\
\caption{Circuito implementado en objetivo particular N°1}
\label{Figura 1}
\end{figure}

\subsection{Actividad N°3}
 Como se puede observar en la imagen anterior, la forma de onda en el canal A es igual a la forma de onda en el canal B, esto se debe a que la linea de transmisión se encuentra adaptada.
 El retardo de tiempo que observamos se debe a la velocidad de propagación de la onda en la linea.

\subsection{Actividad N°4}
En el terminal \textbf{B} de la línea de transmisión, conectar el potenciómetro (ZL= variable).Variar el potenciómetro desde cero al máximo.
Encontrar la posición en que desaparecen las distorsiones, observando la señal a la salida del generador y en la carga. Mantener la posición obtenida por el potenciómetro. 
Cambiar de posición el selector y medir su valor con un óhmmetro.

 \begin{center}

\begin{figure}[hbt]
\centering
\includegraphics[scale = 0.3]{imagen8.jpg}\\
\caption{Circuito implementado en objetivo particular N°1}
\label{Figura 1}
\end{figure}

 
 \textbf{Rpot}=37,71$\Omega$
\end{center}
             

\subsection{Actividad N°5}
Al mover el potenciómetro, variamos la impedancia de la carga y con la ayuda de el osciloscopio pudimos comparar el canal A y B, viendo en la forma de onda la adaptación de la linea

\subsection{Actividad N°6}
6) Con la línea adaptada (ZL = Zo = 75 $\Omega$), medir el tiempo de retardo (\textbf{tr} entre los puntos \textbf{A} y \textbf{B}.
El tiempo de retardo puede ser observado en la primera imagen del trabajo practico

\begin{center}
$t_r$=2 x 0,6 \hspace{0.5cm} $t_r$=1,2$\micro$S
\end{center}

\subsection{Actividad N°7}
 Con la línea adaptada (ZL = Zo = 75 $\Omega$), calcular la atenuación de la línea de transmisión (\textbf{At}) en decibeles (\textbf{dB}).

 
 \begin{center}
\textbf{$E_a$}= 4 $V_{pp}$

\textbf{$E_b$}= 3,2 $V_{pp}$
 
 \vspace{0.5cm}
 
   $A_t$= 20.log  $(\frac{\textbf{$E_b$}}{\textbf{$E_a$}})$  \hspace{0.5cm}   \textbf{$A_t$}= 20.log  $(\frac{3,2}{4})$
 
  \vspace{0.5cm}
 
 \textbf{$A_t$}= -1,938 dB
 \end{center}
 
\subsection{Actividad N°8}
Calcular la constante de propagación (\textbf{K = Vp / C}) de la línea de transmisión por intermedio de la velocidad de propagación \textbf{Vp}.
 \begin{center}
     $V_p$=$(\frac{Longitud de la linea de transmision }{t_r})$       \hspace{0.5cm}    $V_p$=$(\frac{305m}{1,2 {\micro}S})$
      
      \vspace{0.5cm}
    
    $V_p$=254166666,7$(\frac{m}{s})$

 \end{center}

\subsection{Actividad N°9}
 En el punto \textbf{B} de línea de transmisión, desconectar la impedancia de carga ZL = 75 $\Omega$; manteniendo el circuito abierto  (ZL = ∞), observar las formas de onda en los puntos \textbf{A} y \textbf{B}.

\subsection{Actividad N°10}
 Graficar la forma de onda de los canales Ch1 y Ch2.

 \begin{figure}[hbt]
\centering
\includegraphics[scale = 0.3]{imagen8.jpg}\\
\caption{Circuito implementado en objetivo particular N°1}
\label{Figura 1}
\end{figure}

\subsection{Actividad N°11}
 Al no haber ipedancia de carga, el circuito esta abierto. Esto produce que gran parte de la onda sea reflejada, superponiendose a la onda emitida y la reflejada.

\subsection{Actividad N°12}
 En el punto \textbf{B} de línea de transmisión, cortocircuitar la impedancia de carga (ZL= 0); observar las formas de onda en los puntos \textbf{A} y \textbf{B}.
 
\subsection{Actividad N°13}
 Graficar la forma de onda de los canales Ch1 y Ch2.
 \begin{figure}[hbt]
\centering
\includegraphics[scale = 0.3]{imagen1.jpg}\\
\caption{Circuito implementado en objetivo particular N°1}
\label{Figura 1}
\end{figure}

\subsection{Actividad N°14}
 Al acortocircuitar la carga, la lectura del canal B es practicamente constante y la del canal A, la onda cuadrada se ve modificada.
\subsection{Actividad N°15}
 Graficar la forma de onda de los canales Ch1 y Ch2.
 
\subsection{Actividad N°16}
 Calcular el coeficiente de reflexión de tensión en la carga para los tres casos de impedancia de carga analizados (ZL1=75 $\Omega$; ZL= ∞; ZL3 = 0).
 
$\Gamma_E = \frac{Z_n-1}{Z_n+1}$  \vspace{0.5cm}

$Z_n = {\frac{Z_L}{Z_0}}$    \vspace{0.5cm}

\underline{Para $Z_{L1}=75\Omega$}
\begin{center}
$Z_n = {\frac{75 \Omega}{75 \Omega}}$ \hspace{0.5cm}   $Z_n = 1$ \hspace{0.5cm} 
\vspace{0.5cm}

$\Gamma_E = \frac{1-1}{1+1}$  \hspace{0.5cm} $\Gamma_E = 0$ 
\end{center}

\underline{Para $Z_{L2}=\infty$}
\begin{center}
$Z_n = {\frac{\infty}{75 \Omega}}$ \hspace{0.5cm}  $Z_n = \infty$ \hspace{0.5cm} 
\vspace{0.5cm}

$\Gamma_E = \frac{\infty-1}{\infty+1}$  \hspace{0.5cm} $\Gamma_E = 1$ 
\end{center}


\underline{Para $Z_{L3}=0\Omega$}
\begin{center}
$Z_n = {\frac{0\Omega}{75 \Omega}}$ \hspace{0.5cm}  $Z_n = 0$ \hspace{0.5cm} 
\vspace{0.5cm}

$\Gamma_E = \frac{0-1}{0+1}$  \hspace{0.5cm} $\Gamma_E = -1$ 
\end{center}

%\begin{figure}[H]
%\centering
%\includegraphics[scale = 0.5]{gra1.jpg}\\
%\caption{Gráfico de los datos relevados en la tabla N°1}
%\label{Figura 2}
%\end{figure}


\end{document}